\documentclass{experimentreport}

\input{listing_style.tex}

%========================
% 基本信息
%========================

\setcoursename{基础软件开发技术实践}
\setreportname{智能软件工程实验报告}

\setexperimenttitle{\hintholder{请输入实验标题}}

\setstudentname{\hintholder{刘明卓}}
\setdepartment{\hintholder{计算学部}}
\setstudentid{\hintholder{2023111648}}

\setcourseteacher{刘佳琨}
\setinstructor{刘佳琨}

\setlablocation{\hintholder{格物楼213}}
\setlabtime{    \hintholder{7.10}}

\begin{document}

%========================
% 自动生成封面
%========================

\makereporthead

%========================
% 实验目的
%========================

\experimentcontent{
\begin{enumerate}
    \item 理解AI生成代码与人类编写代码之间的差异。
    \item 掌握AI生成代码数据集的构建与筛选方法。
    \item 能够利用实验二、实验三和实验四构建的代码审查模型分析AI生成代码。
    \item 掌握不同代码审查模型在AI生成代码上的性能评估方法。
    \item 理解AI生成代码给代码审查带来的新挑战。
    \item 能够分析不同模型在AI生成代码场景下的优缺点。
\end{enumerate}
}

%========================
% 实验内容
%========================

\experimentcontent{
\begin{enumerate}
    \item 从实验一构建的数据集中筛选AI生成代码对应的PullRequest；
    \item 利用实验二构建的传统机器学习模型进行代码审查；
    \item 利用实验三构建的深度学习模型进行代码审查；
    \item 利用实验四构建的大语言模型进行代码审查；
    \item 比较不同模型在人类代码和AI生成代码上的性能变化；
    \item 分析AI生成代码给代码审查带来的新挑战。
\end{enumerate}
}

%========================
% 实验结果
%========================

\experimentresult{

\placeholder{5cm}{请填写实验结果}

}

%========================
% 问题总结
%========================

\experimentproblems{

\placeholder{5cm}{请总结实验中遇到的问题及解决方法}

}

%========================
% 实验体会
%========================

\experimentexperience{

\placeholder{5cm}{请分享实验体会和收获}

}

%========================
% 教师评语
%========================

\teachercomment

\end{document}